\documentclass[10pt,a4paper]{amsart}
\usepackage[margin=19mm]{geometry}
\usepackage{amsmath,amssymb,mathtools}
\usepackage[scr=rsfs,cal=euler]{mathalfa}
\usepackage{xfrac}
\usepackage[unicode,hidelinks,bookmarks=false]{hyperref}
\newcommand{\ie}{i.e.\ }
\newcommand{\cf}{cf.\ }
\newcommand{\resp}{respectively\ }
\newcommand{\Subsection}{\subsection}
\providecommand{\nobreakdash}{\nobreak\hbox{-}}
\setlength{\emergencystretch}{2em}
\multlinegap0pt
\usepackage{bm}
\usepackage{bbm}
\usepackage{dsfont}
\usepackage{xspace}
\usepackage{cancel}
\usepackage{mathtools}
\usepackage{amsthm}
\makeatletter
\@ifundefined{corollary}{\newtheorem{corollary}{Corollary}}{}
\@ifundefined{lemma}{\newtheorem{lemma}{Lemma}}{}
\@ifundefined{proposition}{\newtheorem{proposition}{Proposition}}{}
\@ifundefined{remark}{\newtheorem{remark}{Remark}}{}
\makeatother
\newcommand{\Ebanach}{{\textnormal E}}
\newcommand{\Hilb}{{\textnormal H}}
\newcommand{\PS}{(\textnormal{PS})^{*}_\tau}
\newcommand{\R}{\mathbb{R}}
\newcommand{\X}{{\textnormal X}}
\newcommand{\ball}{B_\boK}
\newcommand{\boC}{\mathcal{C}}
\newcommand{\boI}{\mathcal{I}}
\newcommand{\boK}{\mathcal{K}}
\newcommand{\minq}{{\bf q}}
\newcommand{\taumax}{T}
\title[A Local Linking Theorem for Relativistic Action Functionals]{A Local Linking Theorem for Relativistic Action Functionals}
\author{Manuel Garz\'on, Salvador L\'opez-Mart\'inez}
\date{Source: arXiv:2607.01845v1 (CC BY 4.0)}
\begin{document}
\maketitle

\noindent\emph{Source and licence.} A Local Linking Theorem for Relativistic Action Functionals. Version arXiv:2607.01845v1 (CC BY 4.0), distributed under the Creative Commons Attribution 4.0 International licence, as recorded in the arXiv licence metadata for this exact version and shown on the abstract page https://arxiv.org/abs/2607.01845v1. Full rights notice: licenses/arxiv-2607.01845-CC-BY-4.0.txt.\par\smallskip

\noindent\textit{Editorial scope.} Counted unit: the Proposition whose source label is \texttt{Pro:Gamma} in the article (source0.tex lines 723--739, source label \texttt{Pro:Gamma}), delivered together with the article's own complete original proof of it (source0.tex lines 740--796, one \texttt{proof} environment of 57 lines). No mathematics is written, repaired or re-derived: statement and proof are the article's own text line for line. Required context retained uncounted: eq:embedding (lines 175--177); def:ball (lines 184--186); hyp:convex (lines 273--275); lemma-alpha (lines 378--384); lemma-LevelSets (lines 426--432); def:regularizedFunctional (lines 447--449); pro:regularizing (lines 451--459); remark:Lipschitz (lines 467--469); H:negative\_levels (lines 532--534); def:omega (lines 549--551); H:existence\_unique\_minimizer (lines 586--588); H:zero\_isolated (lines 590--592); def:Constant (lines 604--606); ineq:regularized-flow (lines 632--634); def:lambda (lines 662--664); eq:charact-lambda-flow (lines 666--668); cor:continuity (lines 676--681). Every definition, display and label of those lines is retained unchanged. Omitted from this package and not redistributed: 1--174, 178--183, 187--272, 276--377, 385--425, 433--446, 450--450, 460--466, 470--531, 535--548, 552--585, 589--589, 593--603, 607--631, 635--661, 665--665, 669--675, 682--722, 797--1590. Nothing inside a retained excerpt is omitted or altered: every retained body is delivered exactly as the article's lines are written, so no omission receipt of any kind accompanies this package. Citations: the retained excerpts carry no citation command at all, so no bibliography entry of the article is transcribed and no reference is redistributed. Reference adaptation: none was needed and none was made; every reference inside the delivered unit resolves inside it, so no unresolved reference or citation is produced. Printed slips kept literal: no printed word, symbol, hypothesis or number of the counted statement or of its proof was altered. Layout and class: the article's own class is \texttt{article} with packages including geometry, amsmath, amsthm, amscd, amssymb, amsfonts, graphics, xcolor, cancel, comment, enumitem, appendix, hyperref, mathtools; the wrapper is the lane's pinned \texttt{amsart} editorial wrapper at 10pt on A4, which restates the theorem-like environments the retained text needs and adds the article's own macro definitions that the retained text needs (\texttt{\textbackslash Ebanach}, \texttt{\textbackslash Hilb}, \texttt{\textbackslash PS}, \texttt{\textbackslash R}, \texttt{\textbackslash X}, \texttt{\textbackslash ball}, \texttt{\textbackslash boC}, \texttt{\textbackslash boI}, \texttt{\textbackslash boK}, \texttt{\textbackslash minq}, \texttt{\textbackslash taumax}), with extra wrapper packages (bm, bbm, dsfont, xspace, cancel, mathtools). The delivered numbering is the unit's own. Assets: no retained span contains a figure environment, a graphics-inclusion command, a PDF-page inclusion, a table or a TikZ picture; no image file is copied into this package, the package declares no asset dependency and no image file is redistributed with it. Licence: version arXiv:2607.01845v1 of this article is distributed under the Creative Commons Attribution 4.0 International licence, as recorded in the arXiv licence metadata for this exact version and shown on the abstract page https://arxiv.org/abs/2607.01845v1. A full rights notice is delivered at licenses/arxiv-2607.01845-CC-BY-4.0.txt. Characters: every retained excerpt is pure ASCII, so the delivered engine, XeLaTeX with its default Latin Modern text font, reports no missing character.
    \begin{equation}\label{eq:embedding}
\|q\|_{\textnormal E}\leq c_0 \|q\|_\Hilb,\quad\hbox{for all }q\in\Hilb.
\end{equation}

\begin{equation}\label{def:ball}
\ball(q,\eta):=\{p\in\boK:\ \|p-q\|_{\textnormal E}\leq \eta\}
\end{equation}

\begin{equation}\label{hyp:convex}
		\hbox{$\boI + \mu\|\cdot\|_\Hilb^2$}\text{ is convex},
\end{equation}

\begin{lemma}\label{lemma-alpha}
	Let $\boI\in\X(\boK)$. Then, for any $q\in\boK$ and any $\alpha>0$, there exists $\eta>0$ such that
	\begin{equation}\label{ineq:lemma1}
		\boI(tp+(1-t) q)<\boI(p)+\alpha,\quad\hbox{for all }(t,p)\in [0,1]\times\ball(q,\eta),
	\end{equation} 
	where the ball $\ball$ is defined in \eqref{def:ball}.
\end{lemma}

\begin{lemma}\label{lemma-LevelSets}
	Let $\boI\in\X(\boK)$ satisfy $\PS$ and have a unique global minimizer $\minq\in\Hilb$. Then, for all $\eta>0$, there exists $\delta>0$ such that
	\begin{equation}\label{eq:InclusionLevelsets}
		\{q\in\boK:\ \boI(q)\leq\boI(\minq)+\delta\}\subset B_\boK(\minq,\eta),
	\end{equation}
    where $\ball$ is defined in \eqref{def:ball}.
\end{lemma}

\begin{equation}\label{def:regularizedFunctional}
	\boI_\varepsilon(q)=\inf_{p\in\Hilb}\left\{\frac1\varepsilon\|p-q\|^2_{\Hilb} + \boI(p)\right\},\ \hbox{ for all }q\in\Hilb.
\end{equation}

\begin{proposition}\label{pro:regularizing} Let $\boI\in\X(\boK)$ satisfy \eqref{hyp:convex}. Then there exists $\varepsilon_0>0$ such that \eqref{def:regularizedFunctional} is finite-valued and of class $\boC^1(\Hilb;\R)$,  for all $\varepsilon\in(0,\varepsilon_0)$. Moreover, $\boI'_\varepsilon:\Hilb\to\Hilb^*$ is globally Lipschitz, and the following properties hold:
	\begin{enumerate}
		\item[(i)]$ \boI(q)\geq \boI_\varepsilon(q) \geq \inf_{p\in\Hilb} \boI(p)$, for all $q\in\Hilb$. 
		\item[(ii)] The critical points of $\boI$ and $\boI_\varepsilon$ are characterized by
		\begin{equation}\label{def:CriticalPoint-Iepsilon}
			\Sigma(\boI)=\Sigma(\boI_\varepsilon)=\{q\in\Hilb:\ \boI(q)=\boI_\varepsilon(q)\}.
		\end{equation}
	\end{enumerate}
\end{proposition}

\begin{remark}\label{remark:Lipschitz}
By the Riesz representation theorem we may identify $\Hilb^*$ with $\Hilb$, under which Proposition~\ref{pro:regularizing} implies that the map $\boI'_\varepsilon:\Hilb\to\Hilb$ is globally Lipschitz continuous.
\end{remark}

\begin{equation}\label{H:negative_levels}
0\geq\boI(q)>\inf_\Hilb\boI,\ \hbox{for all }q\in\Hilb_1\ \hbox{ with }\|q\|_\Hilb\in[0,r],
\end{equation}

\begin{equation}\label{def:omega}
	\boI(\omega(q))=\boI_\varepsilon(q)-\dfrac{1}{\varepsilon}\|q-\omega(q)\|_\Hilb^2,\quad\hbox{and}\quad   \boI_\varepsilon'(q)=\frac{2}{\varepsilon}(q-\omega(q)),\quad \hbox{for all }q\in\Hilb.
\end{equation}

\begin{equation}\label{H:existence_unique_minimizer}
    \Sigma(\boI)\cap \{q\in\Hilb: \boI(q)<0\}=\{\minq\},\quad\boI(\minq)=\min_\Hilb\boI,
\end{equation}

\begin{equation}\label{H:zero_isolated}
  \Sigma(\boI)\cap\{q\in\Hilb_1: \|q\|_\Hilb\in(0, r]\}=\emptyset.
\end{equation}

\begin{equation}\label{def:Constant}
\sigma\in(0,\sigma_r),\quad\quad\varepsilon\in(0,\varepsilon^*(\sigma)).
\end{equation}

\begin{equation}\label{ineq:regularized-flow}
\boI_\varepsilon(x(t;y))=\boI_\varepsilon(y)-t,\quad \hbox{for all } (t,y)\in[0,\taumax(y))\times S_r,
\end{equation}

\begin{equation}\label{def:lambda}
\lambda_\delta(y)=\boI_\varepsilon(y)-\boI(\minq)-\delta,
\end{equation}

\begin{equation}\label{eq:charact-lambda-flow}
\boI_\varepsilon(x(\lambda_\delta(y);y))=\boI_\varepsilon(y)-\lambda_\delta(y)=\boI(\minq)+\delta,\quad\hbox{for all }y\in S_r.
\end{equation}

\begin{corollary}\label{cor:continuity}
Let $\boI\in\X(\boK)$ satisfy $\PS$, \eqref{hyp:convex}, \eqref{H:negative_levels}, \eqref{H:existence_unique_minimizer}, and \eqref{H:zero_isolated}, for some $\minq\in\boK$ and $r\in(0,\|\minq\|_\Hilb)$. Let $\sigma$ and $\varepsilon$ be given by \eqref{def:Constant}. Fix $\delta\in(0,\sigma)$ and define the sets 
\begin{equation}\label{def:LambdaSet}
    \Lambda^r_\delta:=\{(t,y):\  t\in[0,\lambda_\delta(y)]\hbox{ and } y\in S_r\},\quad\Omega^r_\delta:=\{z(t,y):\ (t,y)\in\Lambda^r_\delta\}.
\end{equation} Then, the map $z:\Lambda^r_\delta\to\Omega^r_\delta$ defined by $z(t,y)=x(t;y)$ is Lipschitz. 
\end{corollary}

\begin{proposition}\label{Pro:Gamma}
Let $\boI\in\X(\boK)$ satisfy $\PS$, \eqref{hyp:convex}, \eqref{H:negative_levels}, \eqref{H:existence_unique_minimizer}, \eqref{H:zero_isolated}, for some $\minq\in\boK$ and $r\in(0,\|\minq\|_\Hilb)$. Let $\sigma$ and $\varepsilon$ be given by \eqref{def:Constant}. Then there exists $\delta\in(0,\sigma)$ such that the map $\gamma^*:[0,1]\times S_r\to\boK$, defined by
\begin{equation}\label{def:gamma*}
    \gamma^*(s,y):=\left\{\begin{array}{cl}
         \omega(x(2s\lambda_\delta(y);y)),&\hbox{for }s\in\left[0,\tfrac{1}{2}\right],  \vspace{1mm} \\ 
         (2s-1)\minq+2(1-s)\omega(x(\lambda_\delta(y);y)),&\hbox{for } s\in\left[\tfrac{1}{2},1\right],
    \end{array}\right.
\end{equation}    
is continuous, where $\omega:\Hilb\to\boK$ and $\lambda_\delta:S_r\to\R$ are given by \eqref{def:omega} and \eqref{def:lambda}, respectively. Moreover,
\begin{equation}
\label{eq:gamma-negativeflow}
    \boI(\gamma^*(s,y))\leq0,\quad  \hbox{for all }(s,y)\in[0,1]\times S_r,
\end{equation}and there exists $\rho_\delta\in (0,r)$ satisfying the following:
\begin{equation}\label{eq:gamma-notzero}
    \|\gamma^*(s,y)\|_\Hilb\geq \rho_\delta,\quad \hbox{for all }(s,y)\in[0,1]\times S_r.
\end{equation}
\end{proposition}

\begin{proof}
Let $\delta\in (0,\sigma)$ be chosen small enough later. First, recall from Remark~\ref{remark:Lipschitz} and \eqref{def:omega} that $\omega:\Hilb\to\boK$ is a Lipschitz operator. Thus, the fact that $\gamma^*(s,y)\in\boK$ for every $s\in [0,1/2]$ is obvious. In addition, since $\minq\in\boK$ and $\boK$ is convex, it follows that $\gamma^*(s,y)\in\boK$ for every $s\in (1/2,1]$ as well. Moreover, note that we may write 
\begin{equation}
    \gamma^*(s,y)=\omega(z(2s\lambda_\delta(y),y)),\quad\hbox{for all }(s,y)\in\left[0,\tfrac{1}{2}\right]\times S_r,
\end{equation}where $z:\Lambda_\delta^r\to\Omega_\delta^r$ is introduced in Corollary~\ref{cor:continuity}. This result  implies that the map $\gamma^*$ defined is continuous. For convenience, we maintain the notation $z(t,y)$ instead of $x(t;y)$ in the rest of the proof.



Let us fix $\alpha=-\boI(\minq)-\sigma$ (note that $\alpha>-\min_{S_r}\boI\geq0$), and let $\eta_\sigma>0$ be given by Lemma \ref{lemma-alpha}. We fix also $\eta\in(0,\min\{\eta_\sigma,\|\minq\|_\Ebanach\})$.

Now we prove \eqref{eq:gamma-negativeflow}. From the definition of $\omega$ in \eqref{def:omega}, and by \eqref{ineq:regularized-flow}, we write
\begin{align*}
    \boI(\omega(z(t,y)))&=\boI_\varepsilon(z(t,y))-\dfrac{1}{\varepsilon}\|z(t,y)-\omega(z(t,y))\|_\Hilb^2
    \\
    &=\boI_\varepsilon(y)-t-\dfrac{1}{\varepsilon}\|z(t,y)-\omega(z(t,y))\|_\Hilb^2,\quad\hbox{for all }(t,y)\in\Lambda_\delta^r.
\end{align*}
Then, using that $\boI_\varepsilon|_{S_r}<0=\boI(0)$, we obtain
\begin{equation}\label{ineq:omega-flow}
\boI(\omega(z(t,y)))<0=\boI(0),\quad\hbox{for all }(t,y)\in\Lambda_\delta^r.
\end{equation}
In particular,
\begin{equation} 
    \boI(\gamma^*(s,y))<0, \quad\hbox{for all }(s,y)\in\left[0,\tfrac{1}{2}\right]\times S_r.
\end{equation}

Let us consider the case $(s,y)\in[1/2,1]\times S_r$. By \eqref{def:omega} and \eqref{eq:charact-lambda-flow}, we deduce
\begin{equation}\label{ineq:proofdelta}
    \boI(\omega(z(\lambda_\delta(y),y)))\leq\boI_\varepsilon(z(\lambda_\delta(y),y))=\boI(\minq)+\delta.
\end{equation}
In other words, $\omega(z(\lambda_\delta(y);y)))\in \{q\in\boK: \boI(q)\leq \boI(\minq)+\delta\}$. Then, for $\eta\in (0,\eta_\sigma)$, Lemma \ref{lemma-LevelSets} provides a number $\delta_\eta>0$ such that, fixing $\delta\in (0,\min\{\sigma,\delta_\eta\})$, one has
\begin{equation}\label{ineq:proofeta}
    \|\omega(z(\lambda_\delta(y),y))-\minq\|_\Ebanach<\eta.
\end{equation}Therefore, by taking $t=2(1-s)$ in the notation of Lemma~\ref{lemma-alpha}, we obtain
\begin{equation}\label{ineq:gamma*-negative}
    \boI(\gamma^*(s,y))\leq\boI(\omega(z(\lambda_\delta(y),y)))+\alpha.
\end{equation}
By \eqref{ineq:proofdelta} and the definition of $\alpha$, it follows that
\begin{equation}\label{ineq:gamma*-negative2}
    \boI(\gamma^*(s,y))\leq\boI_\varepsilon(z(\lambda_\delta(y),y))+\alpha=\delta-\sigma<0.
\end{equation}

To conclude, we focus on proving \eqref{eq:gamma-notzero}. By \eqref{ineq:omega-flow}, $
	\omega(z(t,y))\neq0,$ for all $(t,y)\in\Lambda_\delta^r.$ Hence, the continuity of $\omega(z(\cdot))$ in $\Lambda_\delta^r$ implies
\begin{equation}\label{eq:gamma*-positive1}
    m_\delta:=\min_{\left[0,\tfrac{1}{2}\right]\times S_r
}\|\gamma^*(\cdot)\|_\Hilb= \min_{\Lambda_\delta^r}\|\omega(z(\cdot))\|_\Hilb>0.
\end{equation}
Complementarily, when $(s,y)\in\left[\tfrac{1}{2},1\right]\times S_r$, the embedding \eqref{eq:embedding} yields
\begin{equation}
    c_0\|\gamma^*(s,y)\|_\Hilb\geq\|\gamma^*(s,y)\|_{\textnormal E}\geq\|\minq\|_{\textnormal E}-2(1-s)\|\omega(z(\lambda_\delta(y),y))-\minq\|_{\textnormal E},
\end{equation}
where $c_0$ comes from \eqref{eq:embedding}. Then, by \eqref{ineq:proofeta} we conclude that
\begin{equation}\label{eq:gamma*-positive2}
    c_0\|\gamma^*(s,y)\|_\Hilb\geq\|\minq\|_{\textnormal E}-\eta>0,
\end{equation}
and \eqref{eq:gamma-notzero} is satisfied with $\rho_\delta:=c_0^{-1}\min\{\|\minq\|_{\textnormal E}-\eta,m_\delta\}>0$. In particular, note that we may consider $\rho_\delta<r$, which completes the proof.
\end{proof}

\end{document}
